\begin{abstract}
An effective semi-supervised approach for representation learning involves two parallel networks - a trainable network and a momentum network. Two augmented versions of the input images go through these networks to predict their outputs and thereby learn the representation. However, this method often fails to maximize its full potential. We identify two major issues regarding data utilization: (1) the need for large batch sizes, which reduces efficiency in memory-constrained environments, and (2) the exclusion of unlabeled data during downstream tasks. We address these limitations with a two-step approach. First, we incorporate both original and augmented images into the trainable network and add a weighted loss term to enhance learning with the additional input. After pre-training, pseudo labels are generated in a meta-learning setup using two identical teacher-student networks, where the teacher network progressively outperforms the student by receiving feedback from it. Our approach, Representation Learning with Pseudo Labeling (RLPL), improves representation learning by achieving higher accuracy for half the batch size of existing methods. RLPL demonstrates an average performance gain of \textbf{6.12\%} across the STL-10, CIFAR-10, and SVHN datasets, with individual improvements ranging from  \textbf{2.7\%} to \textbf{15.74\%}.
\end{abstract}